\section{Introduction}
\label{sec:introduction}

Let \(L\ge5\) be odd, let \(M=L^m\) with \(m\in\Z_{\ge2}\), and put
\(\eta_k=L^{-k}\), \(k\in\Z_{\ge0}\). We use the photon torus
\(\eta_k\Z^3/\ell\Z^3\) and its chamber sites \((c,g)\), where \(g\)
is one of the \(48\) signed permutation matrices. Spinors have four
components and are antiperiodic in each coordinate. A polymer is a
face-connected union of the partition \(M\)-cubes. Its size parameter
\(Md_M(X)\) is the infimum of the lengths of connected finite nets
meeting every cube of \(X\); it is called small if \(d_M(X)\le L\).
The net convention and its estimates are given in
Appendix~\ref{app:polymer-geometry}.

An activity is a family of Grassmann polynomials in independent
ordinary generators \(\psi,\bar\psi\) and ordered-pair generators
\(\chi,\bar\chi\). The latter represent covariant H\"older differences
only when the fields are evaluated. We work with the space
\(K_k^{x,\C}(\ell)\) defined in
Section~\ref{sec:activities}: its families are balanced, holomorphic
on \(\widetilde R_k\times\mathfrak X_k\), local in the stated
ordinary and pair indices and in the connection bonds, and covariant
under the specified lattice, independent real-gauge and charge
actions. The real subspace \(K_k^x\) satisfies the additional
antilinear identity there. The two generator weights and family norm are
\[
 h_{k,1}=e_k^{-1/4},\qquad h_{k,2}=e_k^{-1/4+\varepsilon},\qquad
 \|\mathscr E\|_k^x
 =\sup_X e^{\kappa d_M(X)}
   \sup_{\widetilde R_k\times\mathfrak X_k}
                 \|\mathscr E(X)\|_{h_k}.
\]
Here \(0<e_k<1\); the weighted coefficient norm and the two
connection domains are specified in
\eqref{eq:rg-activity-coefficient-norm} and
\eqref{eq:connection-domains}. There is no smallness assumption
on the activity in the following linear theorem.

Set
\[
 \bar\varepsilon=\frac1{200},\quad c_X=54,\quad
 \kappa_0=54\log72,\quad K_0=\frac{e^{\kappa_0}}{36},\quad
 c_{\mathrm{rb}}=c_XK_0,\quad c_{\mathrm{fl}}=61c_{\mathrm{rb}}.
\]

\begin{theorem}[Extraction and normalized contraction]\label{thm:main}
Assume that both \(\ell/(LM)\) and \(\ell/(L^2M)\) are odd
integers at least \(13\). Suppose
\begin{equation}\label{eq:rg-contraction-parameters}
 \begin{gathered}
 0<\varepsilon\le\bar\varepsilon,\qquad 1-\varepsilon<\alpha<1,
 \qquad L^{1/8-3\bar\varepsilon/2}\ge2400,\qquad e_k^\varepsilon\le1/4,\\
 160L^{1+3\varepsilon/2}Me_k^\varepsilon\le1,\qquad
 9ML^\varepsilon e_k^{2\varepsilon}\le1,\qquad
 90c_XLMe_k^\varepsilon\le1,
 \end{gathered}
\end{equation}
and
\begin{equation}\label{eq:rg-contraction-polymer-parameters}
 \kappa\ge7\kappa_0+7,\qquad
 L(\kappa-6\kappa_0-4)-1\ge\kappa.
\end{equation}
Let \(\mathscr E\in K_k^{x,\C}(\ell)\) and
\(e_{k+1}=L^{1/2}e_k\).

\begin{enumerate}[label=\textup{(\roman*)},leftmargin=*]
\item There are complex-linear vacuum and mass maps, real on
\(K_k^x\), with
\[
 |\varepsilon_{\mathrm v}(\mathscr E)|
       \le K_0M^{-3}\|\mathscr E\|_k^x,\qquad
 |m(\mathscr E)|
       \le K_0M^{-3}e_k^{1/2}\|\mathscr E\|_k^x .
\]
The extracted mass matrix is \(m(\mathscr E)1_4\).
The local subtraction \(R\) of \eqref{eq:rg-normalized-remainder}
satisfies \(\|R\mathscr E\|_k^x\le2\|\mathscr E\|_k^x\) and the
exact independent-generator identity
\[
 \sum_X\mathscr E(X)
 =-\varepsilon_{\mathrm v}(\mathscr E)\ell^3
       -m(\mathscr E)\langle\bar\psi,\psi\rangle
       +\sum_X(R\mathscr E)(X).
\]
On a small polymer its zero-field vacuum kernel and the integral
of its zero-field ordinary quadratic kernel vanish.

\item The shortest-chain adjustment of
\eqref{eq:rg-adjustment-definition} gives
\(E=(R\mathscr E)^{\mathrm{nat}}\in K_k^{x,\C}(\ell)\), with
\(\|E\|_k^x\le20\|\mathscr E\|_k^x\). It leaves the evaluated
value of \(R\mathscr E\) unchanged whenever
\(\chi=\delta_{\alpha,\mathscr A+\mathscr C}\psi\) and
\(\bar\chi=\delta_{\alpha,-(\mathscr A+\mathscr C)}\bar\psi\).
On every small polymer it has the stronger normalization
\[
 E_{00}(X;0,0)=0,\qquad E_{20}(X;0,0;\cdot,\cdot)=0
\]
as full antisymmetric kernels. Both subtraction and adjustment
leave large-polymer activities unchanged.

\item Scale connections by
\(A_L(b)=L^{-1/2}A(b/L)\), ordinary generators by \(L^{-1}\),
and pair generators by \(L^{-1-\alpha}\).
Let \(\bar X^L\) be the union of \(LM\)-cubes containing \(X\), and set
\[
 (\mathcal LE)(Y;\mathscr A,\mathscr C)
       =\sum_{\bar X^L=LY}S_LE(X;\mathscr A_L,\mathscr C_L).
\]
Then \(\mathscr E\mapsto\mathcal L(R\mathscr E)^{\mathrm{nat}}\)
is a linear map into \(K_{k+1}^{x,\C}(\ell/L)\) and
\[
 \|\mathcal LE\|_{k+1}^x
       \le c_{\mathrm{fl}}L^{-1/4+3\varepsilon/2}\|\mathscr E\|_k^x.
\]
For \(w_j^{\mathrm E}=e_j^{1/4-7\varepsilon}\), the relative estimate is
\[
 \frac{\|\mathcal LE\|_{k+1}^x}{w_{k+1}^{\mathrm E}}
 \le c_{\mathrm{fl}}L^{-3/8+11\varepsilon/2}
                 \frac{\|\mathscr E\|_k^x}{w_k^{\mathrm E}}.
\]
Its factor is at most \(1/16\) if, in addition,
\(c_{\mathrm{fl}}L^{-3/8+6\bar\varepsilon}\le1/16\).
\end{enumerate}
All the stated maps preserve the real subspaces and commute with
setting the second connection argument to zero.
\end{theorem}

The proof reproduces the extraction, adjustment and contraction
arguments of \cite[\S12]{QED3Manuscript}. The first obstruction is
the ordinary quadratic term. Subtracting a local mass makes its
integral vanish, but does not make its kernel vanish at each pair
of points. We transport both legs along shortest chains of cubes
and average their terminal points. The remaining ordinary
quadratic term then factors through the integral of its kernel.
The other terms contain difference generators, whose scaling is
stronger. This construction also explains why the independent
generators are retained until evaluation.

The two lowest degrees require more than degree scaling. On a
small polymer, a real quadratic gauge removes the real constant
and symmetric linear parts of a scaled connection. Two cutoffs,
one for the Taylor remainder and one for the residual linear
field, produce global representatives in an analytic disk of
radius \(L^{7/4}\). The constant complex fields left over are
removed by holomorphic continuation of real-gauge identities.
The zero vacuum value and charge evenness give a second-order
Cauchy gain; the full quadratic normalization gives a first-order
gain. After reblocking, these estimates yield Theorem~\ref{thm:main}.
Sections~\ref{sec:extraction}--\ref{subsec:rg-activity-contraction}
give the proof.

Local reconstruction is a separate estimate, used when unit-scale
mean-field generators are replaced by ordinary and boundary-pair
legs. Theorem~\ref{thm:reconstruction}, stated with its own
hypotheses in Section~\ref{subsec:rg-field-reconstruction}, gives
uniform row bounds for that substitution. The apparent
\(\eta_k^{-1}\) loss in the Dirac stencil cancels before estimation.
The remaining phase defect is summed according to the first scale
at which two neighboring sites have the same block center. Recovery
of a mean field additionally assumes the fine operator is invertible
at the connection in question. No inverse is needed for the row
estimates or for Theorem~\ref{thm:main}.
